\documentclass[10pt,a4paper]{amsart}
\usepackage[margin=19mm]{geometry}
\usepackage{amsmath,amssymb,mathtools}
\usepackage[scr=rsfs,cal=euler]{mathalfa}
\usepackage{xfrac}
\usepackage[unicode,hidelinks,bookmarks=false]{hyperref}
\newcommand{\ie}{i.e.\ }
\newcommand{\cf}{cf.\ }
\newcommand{\resp}{respectively\ }
\newcommand{\Subsection}{\subsection}
\providecommand{\nobreakdash}{\nobreak\hbox{-}}
\setlength{\emergencystretch}{2em}
\multlinegap0pt
\usepackage{bm}
\usepackage{bbm}
\usepackage{dsfont}
\usepackage{xspace}
\usepackage{cancel}
\usepackage{mathtools}
\usepackage{amsthm}
\makeatletter
\@ifundefined{claim}{\newtheorem{claim}{Claim}}{}
\@ifundefined{lemma}{\newtheorem{lemma}{Lemma}}{}
\@ifundefined{theorem}{\newtheorem{theorem}{Theorem}}{}
\makeatother
\title[Gromov Hyperbolicity of Substitution graphs]{Gromov Hyperbolicity of Substitution graphs}
\author{Qingcheng Zeng, Cheng Zeng, Yumei Xue, Huixia He}
\date{Source: arXiv:2608.22836v1 (CC BY 4.0)}
\begin{document}
\maketitle

\noindent\emph{Source and licence.} Gromov Hyperbolicity of Substitution graphs. Version arXiv:2608.22836v1 (CC BY 4.0), distributed under the Creative Commons Attribution 4.0 International licence, as recorded in the arXiv licence metadata for this exact version and shown on the abstract page https://arxiv.org/abs/2608.22836v1. Full rights notice: licenses/arxiv-2608.22836-CC-BY-4.0.txt.\par\smallskip

\noindent\textit{Editorial scope.} Counted unit: the Lemma whose source label is \texttt{le-kakb} in the article (source0.tex lines 1698--1710, source label \texttt{le-kakb}), delivered together with the article's own complete original proof of it (source0.tex lines 1712--1802, one \texttt{proof} environment of 91 lines). No mathematics is written, repaired or re-derived: statement and proof are the article's own text line for line. Required context retained uncounted: th-lac (lines 1479--1493); clm0 (lines 1509--1516); le-GPP (lines 1565--1568); never (lines 1614--1616); clm1 (lines 1626--1630); clm2 (lines 1643--1647). Every definition, display and label of those lines is retained unchanged. Omitted from this package and not redistributed: 1--1478, 1494--1508, 1517--1564, 1569--1613, 1617--1625, 1631--1642, 1648--1697, 1711--1711, 1803--2086. Nothing inside a retained excerpt is omitted or altered: every retained body is delivered exactly as the article's lines are written, so no omission receipt of any kind accompanies this package. Citations: the retained excerpts carry no citation command at all, so no bibliography entry of the article is transcribed and no reference is redistributed. Reference adaptation: none was needed and none was made; every reference inside the delivered unit resolves inside it, so no unresolved reference or citation is produced. Printed slips kept literal: no printed word, symbol, hypothesis or number of the counted statement or of its proof was altered. Layout and class: the article's own class is \texttt{amsart} with packages including amsmath, amssymb, amsfonts, amsaddr, graphicx, booktabs, subcaption, arydshln, color, tikz, cite; the wrapper is the lane's pinned \texttt{amsart} editorial wrapper at 10pt on A4, which restates the theorem-like environments the retained text needs and adds the article's own macro definitions that the retained text needs (none), with extra wrapper packages (bm, bbm, dsfont, xspace, cancel, mathtools). The delivered numbering is the unit's own. Assets: no retained span contains a figure environment, a graphics-inclusion command, a PDF-page inclusion, a table or a TikZ picture; no image file is copied into this package, the package declares no asset dependency and no image file is redistributed with it. Licence: version arXiv:2608.22836v1 of this article is distributed under the Creative Commons Attribution 4.0 International licence, as recorded in the arXiv licence metadata for this exact version and shown on the abstract page https://arxiv.org/abs/2608.22836v1. A full rights notice is delivered at licenses/arxiv-2608.22836-CC-BY-4.0.txt. Characters: every retained excerpt is pure ASCII, so the delivered engine, XeLaTeX with its default Latin Modern text font, reports no missing character.
\begin{theorem}
\label{th-lac}Suppose $K$ is nilpotent and the nilpotency index of $K$ is $I$%
. Let 
\begin{equation*}
S_{1}=\{j\in \Sigma :\text{the }j\text{-th column of }K\text{ is not the
zero vector}\}
\end{equation*}%
and 
\begin{equation*}
S_{2}=\{i\in \Sigma :\text{the }i\text{-th row of }K\text{ is not the zero
vector}\}.
\end{equation*}%
If $\overrightarrow{G}$ satisfies the $(S_{1},S_{2},I-1)$-LAC and $\mathcal{G%
}$ is acyclic, then $\mathcal{U}$ is hyperbolic.
\end{theorem}

\begin{claim}
\label{clm0}Suppose $\overrightarrow{G}$ satisfies the $(S_{1},S_{2},I-1)$%
-LAC and $(u,v),(v,w)\in E_{h}$. If $(u,v)$ is of $J$-type and $(v,w)$ is of 
$G$-type, then the two edges have the same direction$,$ i.e. 
\begin{equation*}
\mathfrak{d}(u,v)=\mathfrak{d}(v,w).
\end{equation*}
\end{claim}

\begin{lemma}
\label{le-GPP}If $\mathcal{G}$ is acyclic, then $\mathcal{U}$ satisfies the
GPP.
\end{lemma}

\begin{equation}
x_{k}^{-}=x_{k+1}^{-}=y_{i}\quad \text{never holds}.  \label{never}
\end{equation}

\begin{claim}
\label{clm1}Suppose $K$ is nilpotent and the nilpotency index of $K$ is $I$.
Let $p$ be a shortest horizontal path in $\mathcal{U}$ with $L(p)>L_{0}$. If 
$p$ falls into Case A, then $M(p)=\boldsymbol{0}.$
\end{claim}

\begin{claim}
\label{clm2}Let $p\ $be a shortest horizontal path in $\mathcal{U}$ with $%
L(p)>L_{0}$. Under the same assumptions as in Theorem \ref{th-lac}, if $p$
falls into Case B, then $M(p)=\boldsymbol{0}.$
\end{claim}

\begin{lemma}
\label{le-kakb}Let $p$ be a shortest horizontal path in $\mathcal{U}$ with $%
L(p)>L_{0}$. Suppose the assumptions of Theorem\ \ref{th-lac} hold and $%
M(p)\neq \mathbf{0}$. Then there exist a shortest horizontal path $p^{\prime
}$ in $\mathcal{U}$, letters $a,b\in \{0,1\}$ and integers $k_{a},k_{b}\in 
\mathbb{N}\cup \{0\}$ such that 
\begin{equation*}
\mathfrak{d(}p)=[a]_{k_{a}}\mathfrak{d(}p^{\prime })[b]_{k_{b}},
\end{equation*}%
where the directed sequence $\mathfrak{d(}p^{\prime })$ starts with $a$ and
ends with $b$, $L(p^{\prime })\leq L_{0}$ and $M(p^{\prime })\neq 
\boldsymbol{0}$.
\end{lemma}

\begin{proof}
Claims \ref{clm1} and \ref{clm2} implies that $M(p)=\boldsymbol{0}$ whenever 
$p$ falls into Case A or Case B. Hence it suffices to consider Case C.\
Consequently, internal stay only occur within the $J_{1}(y_{0})$ and $%
J_{1}(y_{\ell }).$ Let $u_{0}$ denote the last child of $y_{0}$ visited by $%
p $ and $u_{\ell }$ denoe the first child of $y_{\ell }$ visited by $p$,
i.e.,%
\begin{equation*}
u_{0}=x_{i_{0}},\text{ }i_{0}=\max \{j:x_{j}^{-}=y_{0}\},\text{ }
\end{equation*}%
\begin{equation*}
u_{\ell }=x_{i_{\ell }},\text{ }i_{\ell }=\min \{j:x_{j}^{-}=y_{\ell }\}.
\end{equation*}%
Let $H_{0}=G[U_{I-1}(\varphi (u_{0}))],$ $H_{\ell }=G[U_{I-1}(\varphi
(u_{\ell }))],$ and set 
\begin{equation*}
a=\mathfrak{d(}y_{0},y_{1})=\mathfrak{d(}x_{i_{0}},x_{i_{0}+1}),
\end{equation*}%
\begin{equation*}
b=\mathfrak{d(}y_{\ell -1},y_{\ell })=\mathfrak{d(}x_{i_{\ell
}-1},x_{i_{\ell }}).
\end{equation*}

Note that $\varphi (u_{0}),\varphi (v_{\ell })\in S_{1}\cup S_{2}$, hence
the induced subgraphs $H_{0}$ and $H_{\ell }$ are (anti-)arborescences
rooted at $\varphi (u_{0})$ and $\varphi (v_{\ell })$, respectively. In
particular, 
\begin{equation*}
i_{0}<h(H_{0}),\ t-i_{\ell }<h(H_{\ell }).
\end{equation*}%
(Otherwise, if $i_{0}\geq h(H_{0})\ $or$\ t-i_{\ell }\geq h(H_{\ell }),$
then by Claim \ref{clm0} we have $S(p)\geq I.$ This implies $M(p)=0,$
contradicting the assumption that $M(p)\neq 0.$) Then we obtain that%
\begin{equation}
\mathfrak{d(}x_{0},x_{1},\cdots ,x_{i_{0}})=\mathfrak{d(}\varphi
(x_{0}),\varphi (x_{1}),\cdots ,\varphi (x_{i_{0}}))=[a]_{i_{0}}  \label{aaa}
\end{equation}%
and 
\begin{equation}
\mathfrak{d}(x_{i_{\ell }},x_{i_{\ell }+1},\cdots ,x_{t})=\mathfrak{d(}%
\varphi (x_{i_{\ell }}),\varphi (x_{i_{\ell }+1}),\cdots ,\varphi
(x_{t}))=[b]_{t-i_{\ell }}.  \label{bbb}
\end{equation}%
Using (\ref{never}) and (\ref{aaa})--(\ref{bbb}) we have 
\begin{equation}
\mathfrak{d(}p)=\mathfrak{d(}x_{0},x_{1},\cdots ,x_{i_{0}})\mathfrak{d(}%
p^{-})\mathfrak{d(}x_{i_{\ell }},x_{i_{\ell }+1},\cdots ,x_{t})=[a]_{i_{0}}%
\mathfrak{d(}p^{-})[b]_{t-i_{\ell }}  \label{decomp}
\end{equation}%
and $M(p^{-})\neq 0$ whenever $M(p)\neq 0$.

Since $\mathcal{G}$ is acyclic, by Lemma \ref{le-GPP} we obtain that $%
\mathcal{U}$ satisfies GPP. Hence $p^{-}$ is also a shortest horizontal
path. We now proceed by induction on $n$.

For the base case $n=2$, note that $p^{-}$ lies at level $1$, i.e. $p^{-}$
lies in $G_{1}\simeq G.$ \newline
If $i_{0}+(t-i_{\ell })<L(p)-L_{0}$, then 
\begin{equation*}
L(p^{-})=L(p)-i_{0}-(t-i_{\ell })>L_{0},
\end{equation*}%
so by the definition of $L_{0}$ we have $M(p^{-})=\boldsymbol{0},$ and hence 
$M(p)=\boldsymbol{0},$ which contradicts $M(p)\neq 0$. \newline
If $i_{0}+(t-i_{\ell })\geq L(p)-L_{0}$, then $L(p^{-})\leq L_{0}$. By (\ref%
{decomp}), we have 
\begin{equation*}
\mathfrak{d(}p)=[a]_{i_{0}}\mathfrak{d(}p^{-})[b]_{t-i_{\ell }}.
\end{equation*}

Assume the statement holds for some $n\geq 2$. Now we consider the shortest
horizontal path $p$ lying at level $n+1$ with $L(p)>L_{0}$ and $M(p)\neq 0$.
If $L(p^{-})\leq L_{0}$, then by (\ref{decomp}), the projected walk $p^{-}$
already provides the desired $p^{\prime }$.

Otherwise $L(p^{-})>L_{0}$. Since $p^{-}$ is a shortest horizontal path
lying at level $n$ and $M(p^{-})\neq 0$, the induction hypothesis yields a
horizontal path $p^{\prime }$ with $L(p^{\prime })\leq L_{0}$ and $%
M(p^{\prime })\neq 0$ such that 
\begin{equation*}
\mathfrak{d(}p^{-})=[a]_{k_{a}}\mathfrak{d(}p^{\prime })[b]_{k_{b}},
\end{equation*}%
where $\mathfrak{d(}p^{\prime })$ starts with $a$ and ends with $b$, and $%
k_{a},k_{b}\in \mathbb{N}\cup \{0\}$. Substituting this into (\ref{decomp})
gives 
\begin{equation*}
\mathfrak{d(}p)=[a]_{i_{0}}\mathfrak{d(}p^{-})[b]_{t-i_{\ell
}}=[a]_{i_{0}}[a]_{k_{a}}\mathfrak{d(}p^{\prime })[b]_{k_{b}}[b]_{t-i_{\ell
}}=[a]_{i_{0}+k_{a}}\mathfrak{d(}p^{\prime })[b]_{t-i_{\ell }+k_{b}},
\end{equation*}%
which is of the required form. This completes the inductive step.
\end{proof}

\end{document}
